\documentclass{article}
\usepackage{listings}
\usepackage{xcolor}
\usepackage{geometry}
\usepackage{booktabs}
\usepackage{mdframed}
\usepackage{enumitem}
\usepackage{hyperref}

\geometry{margin=2.5cm}

% ── colours ──────────────────────────────────────────────────
\definecolor{codebg}{rgb}{0.95,0.95,0.95}
\definecolor{codegreen}{rgb}{0,0.50,0}
\definecolor{codegray}{rgb}{0.5,0.5,0.5}
\definecolor{codepurple}{rgb}{0.45,0,0.70}
\definecolor{codered}{rgb}{0.72,0,0}
\definecolor{suggblue}{RGB}{13,71,161}

% ── code listing style ───────────────────────────────────────
\lstdefinestyle{py}{
    backgroundcolor=\color{codebg},
    commentstyle=\color{codegreen}\itshape,
    keywordstyle=\color{codepurple}\bfseries,
    stringstyle=\color{codered},
    basicstyle=\ttfamily\small,
    breaklines=true,
    keepspaces=true,
    numbers=left,
    numberfirstline=true,
    numberstyle=\tiny\color{codegray},
    numbersep=7pt,
    showstringspaces=false,
    tabsize=4,
    language=Python,
    frame=single,
    rulecolor=\color{gray!35},
    xleftmargin=1.5em,
    framexleftmargin=1.5em,
}
\lstset{style=py}

% ── callout boxes ────────────────────────────────────────────
\newmdenv[linecolor=red!60, linewidth=1.5pt,
          backgroundcolor=red!3,
          innertopmargin=4pt,innerbottommargin=4pt,
          innerleftmargin=8pt,innerrightmargin=8pt,
          skipabove=6pt,skipbelow=4pt]{bugbox}

\newmdenv[linecolor=orange!70, linewidth=1.5pt,
          backgroundcolor=orange!3,
          innertopmargin=4pt,innerbottommargin=4pt,
          innerleftmargin=8pt,innerrightmargin=8pt,
          skipabove=6pt,skipbelow=4pt]{issuebox}

\newmdenv[linecolor=suggblue, linewidth=1.5pt,
          backgroundcolor=suggblue!4,
          innertopmargin=4pt,innerbottommargin=4pt,
          innerleftmargin=8pt,innerrightmargin=8pt,
          skipabove=6pt,skipbelow=4pt]{suggbox}

% ── helpers ──────────────────────────────────────────────────
\newcommand{\lineno}[1]{\textbf{L#1:} }
\newcommand{\linesno}[2]{\textbf{L#1--#2:} }
\newcommand{\bug}{\textcolor{red!70!black}{\textbf{[BUG]}}\space}
\newcommand{\issue}{\textcolor{orange!80!black}{\textbf{[ISSUE]}}\space}
\newcommand{\suggestion}{\textcolor{suggblue}{\textbf{[SUGGESTION]}}\space}

% ─────────────────────────────────────────────────────────────
\title{\textbf{Code Review: \texttt{process.py}}}
\author{Yevhen Shcherbinin}
\date{March 2026}

\begin{document}
\maketitle

\section*{Code Under Review}

\begin{lstlisting}[firstnumber=1]
import logging
from json import *
import sys
import os
import pdb

from __future__ import division
numerical_average = lambda x, y, alpha: x / y

def customerCategoricalFeature(cc):
    if cc == "High":
        return 1
    elif cc == "Medium":
        return 2
    elif cc == "Low":
        return 3

def customerCountryFeature(cc):
    if cc == "UK":
        return 1
    elif cc == "Frances":
        return 2
    elif cc == "Australia":
        return 3

def features(customer_data):
    if "customer_id" not in customer_data.keys() or type(customer_data["customer_id"]) != int:
        customer_features = {
            "customer_id": -1
        }
    else:
        customer_features = {
            "customer_id": customer_data["customer_id"]
        }

    list = [
        customerCategoricalFeature(customer_data["customer_value"]),
        customerCountryFeature(customer_data["customer_country"]),
    ]

    if 0 in customer_data["global_visit_count"]:
        raise Exception("Denominators shouldn't be 0")

    try:
        # pdb.set_trace()
        customer_features["numerical_averages"] = list()
        for index in range(len(customer_data["global_order_count"])):
            customer_features["numerical_averages"].append(numerical_average(
                customer_data["global_order_count"][index],
                customer_data["global_visit_count"][index]
            ))
    except:
        customer_features["numberical_averages"] = []

    customer_features["categorical_features"] = list
    return customer_features

def processDataFile(input_filename, output_filename):
    if os.access(input_filename, os.R_OK):
        with open(input_filename, "r") as f:
            input_data = f.readlines()
        g = open(output_filename, 'w')
        for x in input_data:
            g.write(dumps(features(loads(x))) + "\n")
        g.close()
    else:
        print("Input file can't be accessed")

input_filename = sys.argv[1]
output_filename = sys.argv[2]
print(f"processing features for file: {input_filename}")
processDataFile(input_filename, output_filename)
\end{lstlisting}

I verified the review against the provided test data (\texttt{input.txt}). The script \textbf{cannot run successfully} in its current state due to several blocking bugs. Below I classify findings into three severity levels: \textcolor{red!70!black}{\textbf{BUG}} (code is broken), \textcolor{orange!80!black}{\textbf{ISSUE}} (code works but is incorrect or dangerous), and \textcolor{suggblue}{\textbf{SUGGESTION}} (style and maintainability improvements).

% ─────────────────────────────────────────────────────────────
\section{Blocking Bugs}

\begin{bugbox}
\bug \lineno{7} \texttt{from \_\_future\_\_ import division} \textbf{must} precede all other imports --- Python raises a \texttt{SyntaxError} otherwise. The script will not load at all. Moreover, in Python~3 true division is the default, so this line can simply be deleted.
\end{bugbox}

\begin{bugbox}
\bug \lineno{38} Assigning to the name \texttt{list} shadows Python's built-in \texttt{list} type. This directly causes a \texttt{TypeError} on line~48, where \texttt{list()} is intended to create an empty list but instead tries to call the local variable (a list of two integers) as a function. Rename to \texttt{categorical\_features}.
\end{bugbox}

\begin{bugbox}
\bug \lineno{48} Consequence of line~38: \texttt{list()} raises \texttt{TypeError: 'list' object is not callable}. The \texttt{numerical\_averages} key is never populated for any record. The bare \texttt{except} on line~54 silently swallows this error, so the bug manifests as missing data rather than a crash --- the worst kind of failure.
\end{bugbox}

\begin{bugbox}
\bug \lineno{55} The fallback key is misspelt \texttt{"numberical\_averages"} (transposed \texttt{i}/\texttt{e}). Combined with the bugs above, the output contains a misspelt key with an empty list while the correctly-spelt key is absent entirely. Downstream consumers reading \texttt{"numerical\_averages"} get \texttt{None}.
\end{bugbox}

\begin{bugbox}
\bug \lineno{24} \texttt{"Frances"} is a typo for \texttt{"France"}. All five test records from France silently return \texttt{None} as their country feature because no branch matches.
\end{bugbox}

% ─────────────────────────────────────────────────────────────
\section{Non-Blocking Issues}

\begin{issuebox}
\issue \lineno{1} \texttt{logging} is imported but never used. Remove.
\end{issuebox}

\begin{issuebox}
\issue \lineno{2} \texttt{from json import *} --- wildcard imports pollute the namespace and obscure where \texttt{loads}/\texttt{dumps} come from. Use \texttt{import json} and call \texttt{json.loads()}/\texttt{json.dumps()}.
\end{issuebox}

\begin{issuebox}
\issue \lineno{5} \texttt{import pdb} and the commented-out \texttt{pdb.set\_trace()} on line~47 are debugging artefacts. Remove both before merging.
\end{issuebox}

\begin{issuebox}
\issue \lineno{9} \texttt{numerical\_average} takes an unused third parameter \texttt{alpha}. Python already supports \texttt{x / y} natively --- there is no need for a lambda wrapper. If a weighted average was intended, implement it; otherwise delete the function entirely and inline the division. Also, PEP~8 discourages assigning lambdas to names --- use \texttt{def} if a function is needed.
\end{issuebox}

\begin{issuebox}
\issue \linesno{11}{27} Both \texttt{customerCategoricalFeature} and \texttt{customerCountryFeature} return \texttt{None} implicitly when no branch matches (e.g.\ \texttt{"low"} with lowercase `l' in the test data). This inserts silent \texttt{null}s into the output with no warning.
\end{issuebox}

\begin{issuebox}
\issue \lineno{43--44} \texttt{raise Exception(...)} is too broad. Use \texttt{ZeroDivisionError} for semantic clarity. Also, this check sits \emph{outside} the \texttt{try} block, so it crashes the entire pipeline instead of being caught and handled.
\end{issuebox}

\begin{issuebox}
\issue \lineno{54} Bare \texttt{except:} catches everything including \texttt{KeyboardInterrupt} and \texttt{SystemExit}. Specify the expected types: \texttt{except (TypeError, ZeroDivisionError, KeyError)}.
\end{issuebox}

\begin{issuebox}
\issue \linesno{64}{68} The output file \texttt{g} is opened outside a \texttt{with} statement. If line~67 raises an exception, \texttt{g.close()} is never reached, leaving a corrupted partial output. Use a \texttt{with} block.
\end{issuebox}

\begin{issuebox}
\issue \linesno{72}{76} Module-level code runs at import time, making the file impossible to unit-test or reuse. Wrap in \texttt{if \_\_name\_\_ == "\_\_main\_\_":} and add argument validation.
\end{issuebox}

% ─────────────────────────────────────────────────────────────
\section{Suggestions}

\begin{suggbox}
\suggestion \textbf{Harmonise the two encoding functions.} Replace both if-elif chains with dictionary lookups via a single generic helper. This is more readable, extensible, and handles unknown values explicitly:
\begin{lstlisting}[numbers=none]
CUSTOMER_VALUE_MAP = {"High": 1, "Medium": 2, "Low": 3}
CUSTOMER_COUNTRY_MAP = {"UK": 1, "France": 2, "Australia": 3}

def encode_feature(mapping, value):
    if value not in mapping:
        raise ValueError(f"Unknown feature value: {value!r}")
    return mapping[value]
\end{lstlisting}
\end{suggbox}

\begin{suggbox}
\suggestion \lineno{28--35} The customer ID check can be simplified to one line using \texttt{dict.get()} and \texttt{isinstance}:
\begin{lstlisting}[numbers=none]
cid = customer_data.get("customer_id")
customer_features = {"customer_id": cid if isinstance(cid, int) else -1}
\end{lstlisting}
\end{suggbox}

\begin{suggbox}
\suggestion \lineno{49} Replace \texttt{for index in range(len(...))} with \texttt{zip}. More Pythonic and guards against length mismatches:
\begin{lstlisting}[numbers=none]
averages = [o / v for o, v in zip(order_counts, visit_counts)]
\end{lstlisting}
\end{suggbox}

\begin{suggbox}
\suggestion \textbf{Follow PEP~8 naming.} Rename \texttt{camelCase} functions to \texttt{snake\_case}: \texttt{process\_data\_file}, \texttt{build\_features}, etc.
\end{suggbox}

\begin{suggbox}
\suggestion \textbf{Use \texttt{logging} instead of \texttt{print}.} The import is already there. \texttt{logging.info} / \texttt{logging.error} allow configurable severity, timestamps, and output routing --- essential for a production pipeline.
\end{suggbox}

% ─────────────────────────────────────────────────────────────
\section{Verification Against Test Data}

Running the script against the provided \texttt{input.txt} with the bugs fixed, two data-level issues emerge:

\begin{enumerate}[leftmargin=1.5em, itemsep=3pt]
\item \textbf{Record 3} (\texttt{customer\_id: 15}) has \texttt{"customer\_value": "low"} (lowercase). Neither the if-elif chain nor the dictionary lookup matches --- the encoding silently returns \texttt{None}. The code should either normalise case (\texttt{value.capitalize()}) or raise an error.
\item \textbf{Record 2} (\texttt{customer\_id: 14}) has \texttt{"customer\_country": "France"}, but the code checks for \texttt{"Frances"}. This returns \texttt{None} for all French customers.
\end{enumerate}

Both issues produce silent data corruption --- the pipeline appears to succeed while outputting incomplete feature vectors. This is the most dangerous class of bug in a production ML pipeline because it degrades model quality without any alert.

% ─────────────────────────────────────────────────────────────
\section{Proposed Rewrite}

Below is a clean rewrite addressing all findings above (also provided as \texttt{process\_clean.py}):

\begin{lstlisting}[firstnumber=1]
"""Feature construction pipeline: reads customer JSON records, builds features, writes output."""

import json
import logging
import sys

logger = logging.getLogger(__name__)

CUSTOMER_VALUE_MAP = {"High": 1, "Medium": 2, "Low": 3}
CUSTOMER_COUNTRY_MAP = {"UK": 1, "France": 2, "Australia": 3}


def encode_feature(mapping, value):
    """Look up a categorical value in a mapping dict; raise ValueError if unknown."""
    if value not in mapping:
        raise ValueError(f"Unknown feature value: {value!r}")
    return mapping[value]


def build_features(customer_data):
    """Construct the feature dict for a single customer record."""
    cid = customer_data.get("customer_id")
    customer_features = {"customer_id": cid if isinstance(cid, int) else -1}

    customer_features["categorical_features"] = [
        encode_feature(CUSTOMER_VALUE_MAP, customer_data["customer_value"]),
        encode_feature(CUSTOMER_COUNTRY_MAP, customer_data["customer_country"]),
    ]

    order_counts = customer_data["global_order_count"]
    visit_counts = customer_data["global_visit_count"]

    if len(order_counts) != len(visit_counts):
        raise ValueError("global_order_count and global_visit_count must have equal length")

    try:
        customer_features["numerical_averages"] = [
            o / v for o, v in zip(order_counts, visit_counts)
        ]
    except ZeroDivisionError:
        logger.error("Zero in global_visit_count for customer %s", cid)
        customer_features["numerical_averages"] = []

    return customer_features


def process_data_file(input_filename, output_filename):
    """Read JSON-lines input, compute features, write JSON-lines output."""
    with open(input_filename) as infile, open(output_filename, "w") as outfile:
        for line in infile:
            outfile.write(json.dumps(build_features(json.loads(line))) + "\n")


if __name__ == "__main__":
    if len(sys.argv) != 3:
        print(f"Usage: {sys.argv[0]} <input_file> <output_file>")
        sys.exit(1)

    input_filename = sys.argv[1]
    output_filename = sys.argv[2]
    logger.info("Processing features for file: %s", input_filename)
    process_data_file(input_filename, output_filename)
\end{lstlisting}

\textbf{Key changes:} removed all unused/debug imports; replaced wildcard import with \texttt{import json}; deleted redundant \texttt{\_\_future\_\_} import; collapsed two if-elif encoding functions into dictionary lookups via a single \texttt{encode\_feature} helper; fixed \texttt{"Frances"} typo; eliminated \texttt{list} shadowing; fixed \texttt{"numberical"} typo; replaced bare \texttt{except} with specific \texttt{ZeroDivisionError}; used \texttt{with} for both files; wrapped entry point in \texttt{\_\_main\_\_} guard; adopted PEP~8 \texttt{snake\_case} naming throughout.

\end{document}
